\documentclass[border=0.2ex,svgnames,tikz]{standalone}
\usepackage{amsmath,mathtools}
\usepackage{fontspec}
\setmainfont{Source Serif 4}
\setsansfont{Source Sans 3}
\setmonofont{Source Code Pro}
\usepackage{ctex}
\setCJKmainfont[BoldFont=Source Han Sans CN]{Source Han Serif CN}

\usetikzlibrary{positioning,calc,arrows.meta}

\begin{document}
\begin{tikzpicture}[
    >=Latex,
    font=\sffamily,
    scale=0.025,
    every node/.style={transform shape=false}
  ]

  % Bar 1
  \begin{scope}[shift={(0,0)}]
    % Main Frame
    \draw (0,0) rectangle (90, 250);

    % OS (220-250) -> y: 0-30
    \draw (0,30) -- (90,30);
    \node at (45, 15) {OS};

    % A (160-220) -> y: 30-90
    % Reserved A (120-160) -> y: 90-130
    % Rect for A+Reserved (120-220) -> y: 30-130
    \draw (0,30) rectangle (90, 130);
    % Dashed line at 160 -> y: 90
    \draw[dashed] (0,90) -- (90,90);
    % Arrow from 160 up to 141 -> y: 90 up to 109
    \draw[->] (45,90) -- (45, 109);
    % Text A at 190 -> y: 60
    \node at (45, 60) {A};

    % B + Reserved B (0-100) -> y: 150-250
    \draw (0,150) rectangle (90, 250);
    % Dashed line at 40 -> y: 210
    \draw[dashed] (0,210) -- (90,210);
    % Arrow from 40 up to 21 -> y: 210 up to 229
    \draw[->] (45,210) -- (45, 229);
    % Text B at 80 -> y: 170
    \node at (45, 170) {B};

    % Label
    \node[below=0.5cm] at (45, 0) {为增长的数据段预留};
  \end{scope}

  % Bar 2
  \begin{scope}[shift={(240,0)}]
    % Main Frame
    \draw (0,0) rectangle (90, 250);

    % OS (220-250) -> y: 0-30
    \draw (0,30) -- (90,30);
    \node at (45, 15) {OS};

    % A Block (120-220) -> y: 30-130
    \draw (0,30) rectangle (90, 130);

    % A Bottom Slice (195-220) -> y: 30-55
    \draw (0,55) -- (90,55);
    \node at (45, 36) {A};

    % A Data Slice (170-195) -> y: 55-80
    \draw[dashed] (0,80) -- (90,80);
    \draw[->] (45,80) -- (45, 94);
    \node at (45, 67.5) {A 数据};

    % A Stack Slice (120-140) -> y: 110-130
    \draw[dashed] (0,110) -- (90,110);
    \draw[->] (45,110) -- (45, 96);
    \node at (45, 120) {A 堆栈};


    % B Block (0-100) -> y: 150-250
    \draw (0,150) rectangle (90, 250);

    % B Bottom Slice (75-100) -> y: 150-175
    \draw (0,175) -- (90,175);
    \node at (45, 156) {B};

    % B Data Slice (50-75) -> y: 175-200
    \draw[dashed] (0,200) -- (90,200);
    \draw[->] (45,200) -- (45, 214);
    \node at (45, 187.5) {B 数据};

    % B Stack Slice (0-20) -> y: 230-250
    \draw[dashed] (0,230) -- (90,230);
    \draw[->] (45,230) -- (45, 216);
    \node at (45, 240) {B 堆栈};

    % Label
    \node[below=0.5cm] at (45, 0) {为增长的数据段、堆栈段预留};
  \end{scope}

\end{tikzpicture}
\end{document}
